\section{The Paradox Documented}

\subsection{Cross-Sectional Evidence}

Sugar supply predicts obesity across countries. Using 2018--2020 country means for women, the bivariate correlation between sugar supply and obesity prevalence is $r = 0.68$ ($p < 0.001$, $N = 37$). Controlling for GDP per capita, urbanization, and vegetable oil supply, the partial correlation remains $r = 0.50$ ($p = 0.002$). Figure~\ref{fig:cross_sectional} shows both relationships.

The association is robust to standard cross-sectional stress tests. Leave-one-out analysis yields partial correlations ranging from 0.42 (dropping Mauritania) to 0.63 (dropping Gambia), with no single country driving the result. Adding subregion fixed effects strengthens the partial correlation to $r = 0.63$ ($p < 0.001$). Adding both subregion and GDP controls gives $r = 0.46$ ($p = 0.004$). Replacing Pearson with Spearman rank correlations preserves the pattern. By every conventional cross-sectional criterion, sugar supply is a strong, stable predictor of obesity.

\begin{figure}[htbp]
\centering
\includegraphics[width=\textwidth]{figures/fig1_cross_sectional.png}
\caption{Cross-sectional association between sugar supply and women's obesity, 2018--2020 country means. Left: bivariate ($r = 0.68$). Right: partial correlation after controlling for GDP, urbanization, and vegetable oil supply ($r = 0.50$, $p = 0.002$).}
\label{fig:cross_sectional}
\end{figure}

\subsection{Within-Country Evidence}

The cross-sectional association vanishes when we follow the same countries over time. Table~\ref{tab:fe_cascade} reports the sugar coefficient under progressively demanding specifications, all with standard errors clustered at the country level \cite{bertrand2004}.

\begin{table}[htbp]
\centering
\caption{Sugar supply coefficient across specifications, controlling for vegetable oil supply. Dependent variable: women's obesity prevalence (\%). Standard errors clustered by country for FE specifications ($N_{\text{clusters}} = 37$, $N_{\text{obs}} = 481$).}
\label{tab:fe_cascade}
\small
\begin{tabular}{lccccc}
\toprule
Specification & $\hat{\beta}_{\text{sugar}}$ & SE & $t$ & $p$ & $R^2$ \\
\midrule
Pooled OLS & 0.630 & 0.032 & 19.84 & $< 0.001$ & 0.49 \\
Country FE & 0.035 & 0.136 & 0.25 & 0.80 & 0.96 \\
Two-way FE & $-0.029$ & 0.045 & $-0.64$ & 0.52 & 0.99 \\
Detrended & $-0.001$ & 0.006 & $-0.13$ & 0.90 & 0.01 \\
\bottomrule
\end{tabular}
\end{table}

Pooled OLS, which mixes between-country and within-country variation, yields a strong sugar coefficient ($t = 19.84$). Adding country fixed effects absorbs the between-country signal and the coefficient drops to $t = 0.25$. Adding year fixed effects drops it further to $t = -0.64$. Entity-specific linear detrending, which removes each country's individual trajectory, yields $t = -0.13$.

Wild-cluster bootstrap $p$-values match the asymptotic clustered results: $p = 0.80$ for country fixed effects and $p = 0.58$ for two-way fixed effects (999 bootstrap draws). A permutation test that shuffles within-country sugar ordering yields $p = 0.58$. Two-way clustering by country and year gives $t = -0.63$, identical to one-way clustering. The null does not depend on the choice of standard error.

Three additional specifications confirm the null. First differences correlate annual changes in sugar with annual changes in obesity: $r = -0.03$ ($p = 0.59$, $N = 444$). Long differences, comparing only the 2010 and 2020 endpoints, yield $r = -0.19$ ($p = 0.25$, $N = 37$). Because long differences do not depend on year-to-year variation, they resist the NCD-RisC smoothing discussed in Section 4.

A Mundlak correlated random effects model \cite{mundlak1978} makes the between-within decomposition explicit. The model includes both time-varying sugar supply and its country mean as regressors, so the coefficients on each cleanly estimate within-country and between-country associations. The within-country sugar coefficient is $t = -0.66$; the between-country coefficient is $t = 3.46$. The entire cross-sectional association comes from between-country variation. Within countries, sugar supply and obesity are unrelated.

Figure~\ref{fig:fe_cascade} visualizes the specification cascade.

\begin{figure}[htbp]
\centering
\includegraphics[width=\textwidth]{figures/fig2_fe_cascade.png}
\caption{Sugar supply coefficient under progressively demanding specifications. The coefficient drops from strongly positive (pooled OLS) to null once within-country variation is isolated.}
\label{fig:fe_cascade}
\end{figure}

\subsection{Generality Beyond Sugar}

The null is not specific to sugar. We tested all 10 FAOSTAT food groups under four specifications: cross-sectional correlation, one-way country fixed effects, two-way fixed effects, and first differences. Six of 10 food groups show significant cross-sectional associations with obesity. Two of 10 remain significant under one-way country fixed effects (meat at $t = 3.05$ and vegetables at $t = 4.13$). Zero of 10 survive two-way fixed effects. Zero of 10 survive first differences. Sugar, meat, milk, and several other food groups look strong in cross-country comparisons, but that signal disappears once shared trends are removed. In this panel, no food-supply variable shows a stable within-country association with obesity.

Recent FAOSTAT-based obesity studies emphasize aggregate energy and nutrient supply rather than individual food groups \cite{siminiuc2025,delgadopertinez2026}. We therefore ran a targeted comparator on grand-total energy, fat, and protein supply. Table~\ref{tab:energy_comparator} shows the same boundary. Total energy is null under two-way fixed effects ($t = -1.32$, $p = 0.19$) and remains null after adding GDP and urban controls ($t = -1.03$, $p = 0.30$). First differences are also null ($t = -0.57$, $p = 0.57$). Fat and protein are null across the same checks. A simple country-linear detrended energy residual is positive ($t = 2.98$, $p = 0.003$). But it weakens after adding year fixed effects ($t = 1.74$, $p = 0.08$), and both lagged and lead energy residuals carry signal (lag $t = 2.98$, lead $t = 2.19$). That pattern is residual co-movement, not directional evidence that national energy supply explains obesity change.

\begin{table}[htbp]
\centering
\caption{Grand-total FAOSTAT comparator against total energy and macronutrients. Entries are country-clustered $t$-statistics for women's obesity prevalence. TWFE + controls includes GDP per capita and urban population share.}
\label{tab:energy_comparator}
\small
\begin{tabular}{lrrrr}
\toprule
Feature & TWFE & TWFE + controls & Detrended & First diff. \\
\midrule
Total energy (+100 kcal/cap/day) & $-1.32$ & $-1.03$ & $2.98^{**}$ & $-0.57$ \\
Total fat (+10 g/cap/day) & $-1.59$ & $-1.44$ & $-0.64$ & $-0.99$ \\
Total protein (+10 g/cap/day) & $-0.78$ & $-0.18$ & $1.50$ & $-0.18$ \\
\bottomrule
\multicolumn{5}{l}{\footnotesize $^{**}p < 0.01$. Detrended removes a country-specific linear trend only.}
\end{tabular}
\end{table}

This uniformity points to a structural explanation. Entity and year fixed effects absorb 99.4\% of obesity variance. The residual 0.6\% is too small and too noisy for any food supply variable to predict, regardless of whether that variable has a genuine causal relationship with obesity at the individual level.

\subsection{Diabetes: The Borderline Case}

Diabetes is the one outcome where the within-country sugar signal approaches significance. Under two-way fixed effects, the women's diabetes sugar coefficient is $t = 1.92$ ($p = 0.055$). A Mundlak decomposition gives a within-country coefficient of $t = 1.99$. This is the closest any specification gets to a conventional threshold in the entire analysis.

The explanation is structural. Entity and year fixed effects absorb 95.6\% of diabetes variance, compared to 99.4\% for obesity. That leaves 4.4\% residual variance -- 7.7 times more room for food supply variables to find something. The borderline TWFE coefficient exploits this extra room. But entity-specific linear detrending, which removes each country's individual diabetes trajectory, yields $t = 0.26$ ($p = 0.80$). First differences yield $r = 0.04$ ($p = 0.41$). Long differences yield $r = 0.24$ ($p = 0.15$). The TWFE signal is residual co-trending that detrending kills. Diabetes shows the same paradox as obesity, with more noise in the residual.

\subsection{GBD Dietary Intake Variables}

In exploratory analysis of 15 dietary intake variables from the Global Burden of Disease 2021 study (fruits, vegetables, legumes, whole grains, nuts, red meat, processed meat, sugar-sweetened beverages, sodium, fiber, calcium, omega-3, polyunsaturated fat, trans fat, and milk), the same pattern held: cross-sectional correlations collapsed under fixed effects. These GBD estimates combine survey data with food supply proxies \cite{gbd2019}, so they are not independent of the FAOSTAT measurement issues we discuss in Section 7. The pattern is consistent with the structural explanation, though we present the GBD results as suggestive rather than definitive.
